%% GENERATED by python/paper/build.py from results/ -- do not edit by hand.
%% Rebuild:  .venv\Scripts\python.exe python\paper\build.py --only Argentina_funcOfRho
%% Source:   results/shocks/universal_match_rho*_commonX.csv
\begin{table}[!htb]
\centering
\begin{threeparttable}
\caption{Pension system reform, year 2010, function of $\rho$}
\label{table:Argentina:funcOfRho}
\begin{tabular}{lccc}
\toprule
\textbf{Scenario} & \textbf{Tax rate} & \textbf{Change in savings rate} & \textbf{Avg. workweek (hours)} \\
\midrule 
Pre-reform & 10.92\% & -- & 42.54 \\[1ex] % row: pre
Post-reform, $\rho$ = 0.5 & 10.63\% & $+0.18$ p.p. & 42.56 \\[1ex] % row: rho@0.5
% Post-reform, $\rho$ = 0.6 & 11.64\% & $-0.12$ p.p. & 42.45 \\[1ex]
% Post-reform, $\rho$ = 0.7 & 12.22\% & $-0.28$ p.p. & 42.38 \\[1ex]
% Post-reform, $\rho$ = 0.8 & 12.59\% & $-0.36$ p.p. & 42.34 \\[1ex]
% Post-reform, $\rho$ = 0.9 & 12.85\% & $-0.41$ p.p. & 42.31 \\[1ex]
Post-reform, $\rho$ = 1.0 & 13.00\% & $-0.43$ p.p. & 42.29 \\[1ex] % row: rho@1.0
% Post-reform, $\rho$ = 1.1 & 13.09\% & $-0.43$ p.p. & 42.28 \\[1ex]
% Post-reform, $\rho$ = 1.2 & 13.15\% & $-0.43$ p.p. & 42.27 \\[1ex]
% Post-reform, $\rho$ = 1.3 & 13.17\% & $-0.42$ p.p. & 42.27 \\[1ex]
% Post-reform, $\rho$ = 1.4 & 13.18\% & $-0.40$ p.p. & 42.27 \\[1ex]
% Post-reform, $\rho$ = 1.5 & 13.17\% & $-0.39$ p.p. & 42.27 \\[1ex]
% Post-reform, $\rho$ = 1.6 & 13.15\% & $-0.37$ p.p. & 42.27 \\[1ex]
% Post-reform, $\rho$ = 1.7 & 13.12\% & $-0.36$ p.p. & 42.27 \\[1ex]
% Post-reform, $\rho$ = 1.8 & 13.10\% & $-0.34$ p.p. & 42.27 \\[1ex]
% Post-reform, $\rho$ = 1.9 & 13.06\% & $-0.33$ p.p. & 42.28 \\[1ex]
Post-reform, $\rho$ = 2.0 & 13.01\% & $-0.31$ p.p. & 42.28 \\[1ex] % row: rho@2.0
\bottomrule
\end{tabular}
\begin{tablenotes}[flushleft]
\footnotesize
\item[] \textit{Note:} The reform permanently shifts $\epsilon = 1-\theta + \theta \eta_1 h_1/h$. Each $\rho$ is separately recalibrated to the same pre-reform targets, so the pre-reform tax rate and workweek are common to every $\rho$. The savings rate is savings relative to GDP; its pre-reform level is not targeted (the capital--output ratio is) and varies with $\rho$, so each post-reform row reports the change against its own $\rho$'s pre-reform level in percentage points.
\end{tablenotes}
\end{threeparttable}
\end{table}
